%% Sample Exam 6 -- Medical Image Analysis (ETH Zurich, 2026)
%%
%% Written in the style of Example_Exam_with_Answers.pdf.
%%
%% Compile the version WITH solutions:      pdflatex "\def\withanswers{}%% Sample Exam 6 -- Medical Image Analysis (ETH Zurich, 2026)
%%
%% Written in the style of Example_Exam_with_Answers.pdf.
%%
%% Compile the version WITH solutions:      pdflatex "\def\withanswers{}%% Sample Exam 6 -- Medical Image Analysis (ETH Zurich, 2026)
%%
%% Written in the style of Example_Exam_with_Answers.pdf.
%%
%% Compile the version WITH solutions:      pdflatex "\def\withanswers{}%% Sample Exam 6 -- Medical Image Analysis (ETH Zurich, 2026)
%%
%% Written in the style of Example_Exam_with_Answers.pdf.
%%
%% Compile the version WITH solutions:      pdflatex "\def\withanswers{}\input{Sample_Exam_6.tex}"
%% Compile the blank version (no answers):  pdflatex Sample_Exam_6.tex
%%
\ifx\withanswers\undefined
  \documentclass[11pt]{exam}
\else
  \documentclass[11pt,answers]{exam}
\fi

\usepackage[margin=2.4cm]{geometry}
\usepackage{amsmath,amssymb}
\usepackage{array}
\usepackage[T1]{fontenc}

% ---------------------------------------------------------------- page style
\pagestyle{foot}
\firstpageheader{}{}{}
\firstpagefooter{}{}{}
\runningheader{}{}{}
\runningfooter{}{}{Page \thepage}

% -------------------------------------------------------------- multiplechoice
% Reproduce the "(A) (B) (C) ..." labelling of the original exam.
\renewcommand{\choicelabel}{$\bigcirc$~(\Alph{choice})}

% ---------------------------------------------------------- true/false macros
% \TFT -> the statement is TRUE, \TFF -> the statement is FALSE.
% In the answer version the correct word is set in capitals and bold.
\newcommand{\TFT}{\ifprintanswers\textbf{TRUE}\hspace{1.1em}False\else True\hspace{1.1em}False\fi}
\newcommand{\TFF}{True\hspace{1.1em}\ifprintanswers\textbf{FALSE}\else False\fi}

% ------------------------------------------------------------- writing space
% Reserves room to write in the blank version; collapses in the answer version.
\newcommand{\answerspace}[1]{\ifprintanswers\else\vspace{#1}\fi}

\begin{document}

% ------------------------------------------------------------------ title page
\begin{center}
  {\LARGE\bfseries Sample Exam}\\[2mm]
  {\Large Medical Image Analysis}\\[6mm]
\end{center}

\vspace{4mm}
\noindent Student name:\hrulefill

\vspace{6mm}
\noindent Immatriculation number:\hrulefill

\vspace{10mm}

\begin{center}
  {\large\bfseries Exam Questions}
\end{center}

\noindent\textbf{Note:} All the questions have equal weight.

\vspace{4mm}

\begin{questions}

% =============================================================== Question 1 ===
\question
In axial 2D acquisitions the slices lie in the transverse plane and are stacked along the
head-to-toe direction. Voxel sizes are given as in-plane (two values) and through-plane (the slice
thickness), and the field-of-view is image size $\times$ voxel size.

\vspace{1mm}
\noindent An axial acquisition uses in-plane voxel size $0.7 \times 0.7$~mm$^2$ with 320 pixels in
both directions, and through-plane voxel size $4.0$~mm with 44 slices.

\begin{parts}

  \part What is the field-of-view in the head-to-toe, left-to-right and anterior-to-posterior
  directions, respectively?

  \begin{solution}
    Answer: $176 \times 224 \times 224$~mm$^3$.

    Head-to-toe is the through-plane direction of an axial acquisition: $44 \times 4.0 = 176$~mm.
    Both in-plane directions give $320 \times 0.7 = 224$~mm.
  \end{solution}
\answerspace{2.2cm}

  \part The anatomy that has to be covered extends 210~mm in the head-to-toe direction. Does this
  protocol cover it? If not, how many slices are needed at the same slice thickness, and by what
  factor does the acquisition time change if it is proportional to the number of slices?

  \begin{solution}
    Answer: \textbf{No} --- the protocol covers only 176~mm of the required 210~mm.
    $\lceil 210 / 4.0 \rceil = \lceil 52.5 \rceil = \mathbf{53}$ slices are needed (53 slices cover
    212~mm). The acquisition time grows by a factor $53/44 \approx \mathbf{1.20}$, i.e.\ about
    20\% longer.

    Note that you must round \emph{up}: 52 slices would cover only 208~mm and miss the anatomy.
  \end{solution}
\answerspace{3.2cm}

  \part Compute the anisotropy ratio of the voxels. This volume is then fed, without resampling,
  to a 3D convolutional network that was trained on $1.0$~mm isotropic data. What goes wrong?

  \begin{solution}
    Answer: the anisotropy ratio is $4.0 / 0.7 \approx \mathbf{5.7}$ --- the voxel is almost six
    times longer through-plane than in-plane.

    A $3 \times 3 \times 3$ kernel therefore covers $2.1 \times 2.1 \times 12.0$~mm$^3$ of real
    tissue instead of the $3 \times 3 \times 3$~mm$^3$ it saw during training. The receptive field
    of every neuron is wildly asymmetric in physical space, so the spatial features the network
    learned simply do not correspond to the same physical structures any more. This is a domain
    shift the model was never trained for, and the fix is to resample to the training voxel size
    first.

    Note that resampling does not create the missing information: upsampling the through-plane
    direction from 4.0~mm to 1.0~mm interpolates three samples out of every one that was really
    measured. The effective resolution in that direction is still 4~mm.
  \end{solution}
\answerspace{4.5cm}

\end{parts}

% =============================================================== Question 2 ===
\question
In atlas-based segmentation the atlas is registered to the image and then used inside the mixture
model. Which of the two posteriors below is the one used by atlas-based segmentation, as opposed to
plain EM with a Gaussian mixture model? $\pi_n$ denotes a mixture weight estimated from the image
itself. Please circle the correct choice.

\vspace{2mm}
\begin{center}
  $\displaystyle p(c(x){=}n \mid I(x)) \propto \mathcal{N}(I(x) \mid \mu_n, \Sigma_n)\ \pi_n$
  \hspace{1.2cm}
  $\displaystyle p(c(x){=}n \mid I(x)) \propto \mathcal{N}(I(x) \mid \mu_n, \Sigma_n)\
  p_{atlas}(c(x){=}n)$
\end{center}
\vspace{2mm}

\begin{solution}
  Answer: Right one.

  The atlas prior $p_{atlas}(c(x){=}n)$ \emph{replaces} the free mixture weight $\pi_n$. Two things
  change as a result. The prior is now spatially varying --- it says something different at every
  voxel, which is exactly the information plain EM lacks --- and it no longer has to be estimated,
  so the class probabilities drop out of the optimisation entirely.
\end{solution}
\answerspace{1.5cm}

% =============================================================== Question 3 ===
\question
Registration problems are classified along two axes: whether the two images come from the
\emph{same or different subjects}, and whether they were acquired with the \emph{same or different
modalities}.

\begin{parts}
  \part Which of the four combinations is the easiest and which is the hardest, and why?

  \begin{solution}
    Answer: easiest is \textbf{intra-subject, intra-modality}; hardest is \textbf{inter-subject,
    inter-modality}.

    The same subject imaged twice with the same sequence shows the same anatomy with comparable
    intensities, so both the transformation and the similarity measure have an easy job. Different
    subjects imaged with different modalities combine two independent difficulties: the anatomies
    genuinely differ, and the intensities are not comparable. The relative difficulty of the two
    middle cases is a subjective judgement. Difficulty also depends on the anatomy --- brain is
    considerably easier than abdomen.
  \end{solution}
\answerspace{3.4cm}

  \part Each axis primarily determines one component of the registration algorithm. Which axis
  determines which, and why?

  \begin{solution}
    Answer: the \textbf{subject} axis determines the \textbf{transformation model}; the
    \textbf{modality} axis determines the \textbf{loss function}.

    Same subject means the same anatomy is present in both images, so a rigid or affine model is
    often enough --- what changed is position, not shape. Different subjects have genuinely
    different anatomy, with no exact correspondence and possibly pathology, so a deformable model
    is required.

    The modality determines whether intensities are comparable. Within a modality, SSD is
    appropriate. Across modalities it fails outright, and you need a statistical metric that is
    invariant to the intensity relationship --- normalised cross-correlation for a linear
    relationship, mutual information for an arbitrary one.
  \end{solution}
\answerspace{3.6cm}
\end{parts}

% =============================================================== Question 4 ===
\question
Please indicate whether the statement is true or false by circling the correct choice in each of
the following.

\begin{parts}
  \part \TFT\ \ The $3 \times 3$ mean filter is separable, $S_3 = \frac{1}{9}[1,1,1] * [1,1,1]^T$,
  which makes it cheaper to apply than a general $3 \times 3$ kernel.

  \part \TFF\ \ A Gaussian has finite support, so a $3 \times 3$ structural element represents a
  Gaussian filter exactly.

  \part \TFT\ \ Binomial filters are built by iterated convolution of $[1,1]$ and have better
  properties in the frequency domain than mean filtering.

  \part \TFF\ \ Increasing $\sigma$ of a Gaussian filter smooths more, but the size of structural
  element required is independent of $\sigma$.
\end{parts}

\begin{solution}
  (a) TRUE --- a separable 2D kernel can be applied as two 1D passes, which costs $O(k)$ per pixel
  instead of $O(k^2)$. The Gaussian and binomial kernels are separable as well.

  (b) FALSE --- the Gaussian has \emph{infinite} support. Any finite structural element is only an
  approximation of it, truncated at some distance from the centre.

  (c) TRUE --- $S_3 = \frac{1}{16}\left[\begin{smallmatrix}1&2&1\\2&4&2\\1&2&1\end{smallmatrix}
  \right]$, obtained from successive convolutions of $[1,1]$, and larger binomial kernels from
  successive convolutions of smaller ones. They are integer filters and behave better in the
  frequency domain than the box filter.

  (d) FALSE --- and (b) is the reason. Since the element only approximates the Gaussian by
  truncating it, a wider Gaussian needs a correspondingly larger element for the approximation to
  remain good. Using a $3 \times 3$ element with a large $\sigma$ silently truncates most of the
  kernel.
\end{solution}

% =============================================================== Question 5 ===
\question
A classification network takes an input image $x \in \mathbb{R}^{64 \times 64}$ with a single
channel and maps it through hidden layers of 256, 128 and 64 units to 16 output classes. Consider
only the \emph{first} link, from the input to the 256-unit layer, and ignore biases throughout.

\begin{parts}
  \part How many weights does this link have if it is fully connected?
  \begin{solution}
    Answer: $64 \times 64 \times 256 = 4\,096 \times 256 = \mathbf{1\,048\,576}$ weights.
  \end{solution}
\answerspace{1.8cm}

  \part How many weights does it have if it is instead a valid $5 \times 5$ convolution producing
  256 output channels from the single input channel?
  \begin{solution}
    Answer: $5 \times 5 \times 1 \times 256 = \mathbf{6\,400}$ weights.
  \end{solution}
\answerspace{1.8cm}

  \part What is the ratio, and what property of the convolutional layer explains it?
  \begin{solution}
    Answer: $1\,048\,576 / 6\,400 = \mathbf{163.84}$, so the fully connected link has roughly 164
    times more parameters.

    The explanation is \textbf{weight sharing}: the same $5 \times 5$ kernel is applied at every
    position, so the parameter count depends only on the kernel size and the channel counts --- and
    not at all on the image size. The fully connected count, by contrast, scales with the number of
    pixels. Doubling the image to $128 \times 128$ would quadruple the fully connected count to over
    four million while leaving the convolutional count at exactly 6\,400.

    Note also that the convolutional layer \emph{lifts} the data to 256 channels at every spatial
    position rather than compressing the whole image into 256 numbers, so it retains far more
    information despite using far fewer parameters.
  \end{solution}
\answerspace{3.4cm}
\end{parts}

% =============================================================== Question 6 ===
\question
In atlas-based segmentation the model is $p(I(x) \mid \theta) = \sum_n \mathcal{N}(I(x) \mid \mu_n,
\Sigma_n)\, p_{atlas}(c(x){=}n)$ with $\theta = \{\mu_n, \Sigma_n\}_{n=1}^N$. Please explain why the
intensity parameters $\mu_n$ and $\Sigma_n$ still have to be estimated for every new MR image, but
could in principle be fixed in advance for CT. Then state what else the atlas gives you compared
with plain EM. You do not need to write down the equations.

\begin{solution}
  Answer: \emph{Why MRI needs them fitted.} MRI has \textbf{no absolute intensity scale}. The number
  stored in a voxel depends on the scanner, the coil, the sequence parameters and the
  reconstruction, so ``the intensity of white matter'' is not a fixed quantity --- it differs
  between two scans of the same person on two machines. The Gaussian parameters therefore have to be
  re-estimated from each image.

  \emph{Why CT could fix them.} CT intensities are expressed in \textbf{Hounsfield units}, which are
  physically calibrated against water and air. A given tissue has a reproducible value across
  scanners, so $\mu_n$ and $\Sigma_n$ can be tabulated in advance and all the parameters fixed.

  \emph{What else the atlas buys.} Two things beyond the accuracy gain. It supplies
  \textbf{spatially varying prior information} --- at each voxel it says what is expected there,
  which is exactly what plain EM has no way of knowing. And it \textbf{removes the label
  ambiguity}: unsupervised EM returns clusters in an arbitrary order with no names attached, whereas
  with an atlas component $n$ \emph{is} grey matter. You know what you asked for, and there is no
  randomness left in which component is which.
\end{solution}
\answerspace{7.5cm}

% =============================================================== Question 7 ===
\question
Which of the following statements about interpolation is \textbf{inaccurate}?

\begin{choices}
  \choice Bilinear interpolation can equivalently be written as the parametric surface
  $I(x_1,x_2) = a + b x_1 + c x_2 + d x_1 x_2$, whose four unknowns are fixed by the intensities at
  the four surrounding pixels.
  \choice The trilinear extension adds the remaining monomials to give eight unknowns, fixed by the
  eight corner voxels of the cell containing the query point.
  \choice Bicubic interpolation has 16 unknowns, determined by the four corner intensities together
  with their first and mixed derivatives --- and since those derivatives must themselves be
  approximated from further pixels, bicubic effectively uses more than four samples.
  \choice Because bilinear interpolation is a weighted average of the surrounding values, applying
  it repeatedly during an iterative registration leaves the effective spatial resolution unchanged.
\end{choices}

\begin{solution}
  Answer: D.

  Averaging is exactly what costs you resolution. Bilinear interpolation smooths the image slightly
  every time it is applied, and during an optimisation loop it is applied thousands of times, so the
  blurring compounds. This is one reason the standard strategy is bilinear \emph{during} the
  optimisation, where speed matters, and bicubic for the \emph{final} resampling, where quality
  matters and the cost is paid once.
\end{solution}

% =============================================================== Question 8 ===
\question
Which of the following statements about the loss functions used for image restoration is
\textbf{incorrect}?

\begin{choices}
  \choice MAE, MSE, NMSE, PSNR and SSIM are all computed directly between the prediction and the
  ground-truth image.
  \choice SSIM compares many local patches of the two images and combines three components ---
  luminance, contrast and structure.
  \choice A perceptual loss passes both images through a previously trained network and compares
  intermediate features; that feature network must have been trained on the same task and the same
  modality as the restoration model, otherwise the comparison is meaningless.
  \choice An adversarial loss is a distributional distance defined over \emph{sets} of samples
  rather than between two paired images.
\end{choices}

\begin{solution}
  Answer: C.

  The feature network does not have to match the task or the modality. The lecture explicitly allows
  a well-established network trained on \emph{natural} images --- VGG is the standard choice --- or a
  network trained on CT images for some entirely different task. What is being borrowed is a
  general-purpose notion of what makes two images look different in terms of contextual content
  rather than pixel values, and that transfers surprisingly well across domains.

  The motivation for all of this is that MSE-type losses produce visibly over-smoothed output;
  combining a perceptual loss with an adversarial one gave the most realistic denoised low-dose CT
  in the referenced work.
\end{solution}

% =============================================================== Question 9 ===
\question
A semi-supervised auxiliary loss enforces consistency under spatial transformations. Two
transformations $\phi_1, \phi_2$ are drawn per image and the loss is
$\mathcal{L}\big(\phi_2 \circ \phi_1^{-1} \circ f(\phi_1 \circ x;\theta),\ f(\phi_2 \circ x;\theta)
\big)$. Which of the following statements is \textbf{incorrect}?

\begin{choices}
  \choice The loss requires no labels, so it can be evaluated on unlabelled images.
  \choice It encodes the principle that if $x \Leftrightarrow y$ then $\phi \circ x \Leftrightarrow
  \phi \circ y$, which is the same principle that underlies data augmentation.
  \choice The inverse $\phi_1^{-1}$ appears so that the two predictions are brought into a common
  spatial frame; without it they would sit in different frames and could not be compared pixel by
  pixel.
  \choice Because the loss is defined on unlabelled data, it replaces the supervised loss rather
  than being added to it.
\end{choices}

\begin{solution}
  Answer: D.

  It is an \emph{auxiliary} loss: the total cost is the supervised loss over the labelled examples
  \emph{plus} the consistency term, and the consistency term is summed over labelled and unlabelled
  examples alike. Dropping the supervised loss would remove the only thing that ties the network's
  output to the actual task --- a network that predicts a constant map everywhere is perfectly
  consistent under any transformation, and perfectly useless.
\end{solution}

% ============================================================== Question 10 ===
\question
When the network parameters are sampled from their posterior rather than optimised to a point
estimate, the sampled parameters can perform noticeably worse on the task than the MAP estimate.
What is the explanation given in the lecture?

\begin{choices}
  \choice In high dimensions the posterior is always unimodal, so every sample collapses onto the
  mean and no diversity is possible.
  \choice For an isotropic Gaussian in $d$ dimensions the expected distance of a sample from the
  mean grows like $\sqrt{d}$, so samples concentrate away from the mode --- and since the likelihood
  $p(D \mid \theta, M)$ \emph{is} the task loss, moving away from the mode means moving away from
  well-performing parameters.
  \choice High-dimensional posteriors cannot be sampled at all, so any apparent sample is an
  artefact of the random number generator.
  \choice In high dimensions the task loss becomes independent of the parameters, so performance is
  essentially random.
  \choice Sampling error decreases with dimension, so the effect disappears entirely for large
  networks.
\end{choices}

\begin{solution}
  Answer: B.

  This is the ``typical set'' phenomenon and it is thoroughly counter-intuitive: in high dimensions
  almost none of the probability mass sits near the mode, even though the mode is the single most
  probable point. Samples live in a thin shell at distance $\approx \sqrt{d}$ from the mean.

  The consequence for uncertainty estimation is uncomfortable. Because $\log p(D \mid \theta, M)$ is
  precisely the training objective, a parameter vector drawn honestly from the posterior can sit in
  a region of poor task performance --- so a faithful posterior sample and a well-performing network
  are not the same thing. Bounded sampling schemes have been proposed to mitigate this, but they do
  not fully solve it.
\end{solution}

% ============================================================== Question 11 ===
\question
A Vision Transformer is applied to images of size $224 \times 224$ with 3 channels. It uses a patch
size of $P = 16$, an embedding dimension of $D = 768$, and 12 attention heads. Recall that the image
is split into $N = HW/P^2$ patches, each flattened and mapped by a patch-embedding matrix
$E \in \mathbb{R}^{(P^2 C) \times D}$, and that a learnable class token is prepended to the sequence.
Circle True or False:

\begin{parts}
  \part \TFT\ \ The number of patch tokens is $N = 196$.

  \part \TFT\ \ The patch-embedding matrix $E$ is $768 \times 768$ and therefore contains
  $589\,824$ weights.

  \part \TFF\ \ After the class token is added, the sequence entering the encoder still has length
  196, because the class token takes the place of one patch.

  \part \TFF\ \ The attention matrix $QK^T$ of a single head has $N \times N$ entries, so its size
  grows linearly with the number of tokens.
\end{parts}

\begin{solution}
  (a) TRUE --- $N = HW/P^2 = (224 \times 224)/16^2 = 50\,176/256 = 196$, i.e.\ a $14 \times 14$ grid
  of patches.

  (b) TRUE --- the flattened patch has $P^2 C = 16^2 \times 3 = 256 \times 3 = 768$ entries, and it
  is mapped to $D = 768$, so $E$ is $768 \times 768 = 589\,824$. (That the two dimensions coincide is
  a property of this particular configuration, not a general rule.)

  (c) FALSE --- the class token is \emph{prepended}, not substituted: the sequence length becomes
  $196 + 1 = 197$. It carries no patch content of its own; it exists so that the final
  representation at that position can be read out and passed to the classification MLP.

  (d) FALSE --- an $N \times N$ matrix has $N^2$ entries, so the cost grows \textbf{quadratically}
  with sequence length. Here that is $197^2 = 38\,809$ attention scores per head, and 12 heads per
  layer. This quadratic scaling is the reason patching exists at all: attending over all
  $224 \times 224 = 50\,176$ pixels individually would need over $2.5$ billion scores per head.
\end{solution}

% ============================================================== Question 12 ===
\question
Which of the following statements on self-supervised learning are correct (multiple answers are
possible)?

\begin{choices}
  \choice Contrastive learning needs labelled examples in order to decide which pairs of images
  should be treated as positive and which as negative.
  \choice A pre-text task is a task that requires only images; the network trained on it is then
  fine-tuned with very few labelled examples on the task you actually care about.
  \choice In the \emph{local} contrastive loss the transformations must either preserve the patch
  correspondences or be undone before the corresponding patches are compared.
  \choice Foundation models in this setting are obtained by supervised training at very large scale
  on carefully curated labelled data sets.
  \choice The InfoNCE loss maximises the similarity between two views of the same image through its
  numerator, and minimises the similarity to the samples of the negative set through its
  denominator.
\end{choices}

\begin{solution}
  Answer: (B), (C) and (E) are correct.

  (A) is wrong, and it is wrong about the whole point of the method. The labels in contrastive
  learning are \emph{surrogate} labels manufactured from the data itself: two random transformations
  of the same image are declared ``close'', everything else ``far''. No human annotation is involved
  --- which is precisely why the approach is useful when labels are scarce. Note also that the
  transformations are not restricted to geometric ones; intensity changes are used deliberately, so
  that the representation becomes invariant to them.

  (D) inverts the defining property. Foundation models here are built by \textbf{self-supervised}
  training at very large scale --- masked auto-encoding and variations of contrastive learning ---
  on transformer architectures. Requiring curated labels at that scale is exactly what these methods
  exist to avoid.
\end{solution}

\end{questions}

\end{document}
"
%% Compile the blank version (no answers):  pdflatex Sample_Exam_6.tex
%%
\ifx\withanswers\undefined
  \documentclass[11pt]{exam}
\else
  \documentclass[11pt,answers]{exam}
\fi

\usepackage[margin=2.4cm]{geometry}
\usepackage{amsmath,amssymb}
\usepackage{array}
\usepackage[T1]{fontenc}

% ---------------------------------------------------------------- page style
\pagestyle{foot}
\firstpageheader{}{}{}
\firstpagefooter{}{}{}
\runningheader{}{}{}
\runningfooter{}{}{Page \thepage}

% -------------------------------------------------------------- multiplechoice
% Reproduce the "(A) (B) (C) ..." labelling of the original exam.
\renewcommand{\choicelabel}{$\bigcirc$~(\Alph{choice})}

% ---------------------------------------------------------- true/false macros
% \TFT -> the statement is TRUE, \TFF -> the statement is FALSE.
% In the answer version the correct word is set in capitals and bold.
\newcommand{\TFT}{\ifprintanswers\textbf{TRUE}\hspace{1.1em}False\else True\hspace{1.1em}False\fi}
\newcommand{\TFF}{True\hspace{1.1em}\ifprintanswers\textbf{FALSE}\else False\fi}

% ------------------------------------------------------------- writing space
% Reserves room to write in the blank version; collapses in the answer version.
\newcommand{\answerspace}[1]{\ifprintanswers\else\vspace{#1}\fi}

\begin{document}

% ------------------------------------------------------------------ title page
\begin{center}
  {\LARGE\bfseries Sample Exam}\\[2mm]
  {\Large Medical Image Analysis}\\[6mm]
\end{center}

\vspace{4mm}
\noindent Student name:\hrulefill

\vspace{6mm}
\noindent Immatriculation number:\hrulefill

\vspace{10mm}

\begin{center}
  {\large\bfseries Exam Questions}
\end{center}

\noindent\textbf{Note:} All the questions have equal weight.

\vspace{4mm}

\begin{questions}

% =============================================================== Question 1 ===
\question
In axial 2D acquisitions the slices lie in the transverse plane and are stacked along the
head-to-toe direction. Voxel sizes are given as in-plane (two values) and through-plane (the slice
thickness), and the field-of-view is image size $\times$ voxel size.

\vspace{1mm}
\noindent An axial acquisition uses in-plane voxel size $0.7 \times 0.7$~mm$^2$ with 320 pixels in
both directions, and through-plane voxel size $4.0$~mm with 44 slices.

\begin{parts}

  \part What is the field-of-view in the head-to-toe, left-to-right and anterior-to-posterior
  directions, respectively?

  \begin{solution}
    Answer: $176 \times 224 \times 224$~mm$^3$.

    Head-to-toe is the through-plane direction of an axial acquisition: $44 \times 4.0 = 176$~mm.
    Both in-plane directions give $320 \times 0.7 = 224$~mm.
  \end{solution}
\answerspace{2.2cm}

  \part The anatomy that has to be covered extends 210~mm in the head-to-toe direction. Does this
  protocol cover it? If not, how many slices are needed at the same slice thickness, and by what
  factor does the acquisition time change if it is proportional to the number of slices?

  \begin{solution}
    Answer: \textbf{No} --- the protocol covers only 176~mm of the required 210~mm.
    $\lceil 210 / 4.0 \rceil = \lceil 52.5 \rceil = \mathbf{53}$ slices are needed (53 slices cover
    212~mm). The acquisition time grows by a factor $53/44 \approx \mathbf{1.20}$, i.e.\ about
    20\% longer.

    Note that you must round \emph{up}: 52 slices would cover only 208~mm and miss the anatomy.
  \end{solution}
\answerspace{3.2cm}

  \part Compute the anisotropy ratio of the voxels. This volume is then fed, without resampling,
  to a 3D convolutional network that was trained on $1.0$~mm isotropic data. What goes wrong?

  \begin{solution}
    Answer: the anisotropy ratio is $4.0 / 0.7 \approx \mathbf{5.7}$ --- the voxel is almost six
    times longer through-plane than in-plane.

    A $3 \times 3 \times 3$ kernel therefore covers $2.1 \times 2.1 \times 12.0$~mm$^3$ of real
    tissue instead of the $3 \times 3 \times 3$~mm$^3$ it saw during training. The receptive field
    of every neuron is wildly asymmetric in physical space, so the spatial features the network
    learned simply do not correspond to the same physical structures any more. This is a domain
    shift the model was never trained for, and the fix is to resample to the training voxel size
    first.

    Note that resampling does not create the missing information: upsampling the through-plane
    direction from 4.0~mm to 1.0~mm interpolates three samples out of every one that was really
    measured. The effective resolution in that direction is still 4~mm.
  \end{solution}
\answerspace{4.5cm}

\end{parts}

% =============================================================== Question 2 ===
\question
In atlas-based segmentation the atlas is registered to the image and then used inside the mixture
model. Which of the two posteriors below is the one used by atlas-based segmentation, as opposed to
plain EM with a Gaussian mixture model? $\pi_n$ denotes a mixture weight estimated from the image
itself. Please circle the correct choice.

\vspace{2mm}
\begin{center}
  $\displaystyle p(c(x){=}n \mid I(x)) \propto \mathcal{N}(I(x) \mid \mu_n, \Sigma_n)\ \pi_n$
  \hspace{1.2cm}
  $\displaystyle p(c(x){=}n \mid I(x)) \propto \mathcal{N}(I(x) \mid \mu_n, \Sigma_n)\
  p_{atlas}(c(x){=}n)$
\end{center}
\vspace{2mm}

\begin{solution}
  Answer: Right one.

  The atlas prior $p_{atlas}(c(x){=}n)$ \emph{replaces} the free mixture weight $\pi_n$. Two things
  change as a result. The prior is now spatially varying --- it says something different at every
  voxel, which is exactly the information plain EM lacks --- and it no longer has to be estimated,
  so the class probabilities drop out of the optimisation entirely.
\end{solution}
\answerspace{1.5cm}

% =============================================================== Question 3 ===
\question
Registration problems are classified along two axes: whether the two images come from the
\emph{same or different subjects}, and whether they were acquired with the \emph{same or different
modalities}.

\begin{parts}
  \part Which of the four combinations is the easiest and which is the hardest, and why?

  \begin{solution}
    Answer: easiest is \textbf{intra-subject, intra-modality}; hardest is \textbf{inter-subject,
    inter-modality}.

    The same subject imaged twice with the same sequence shows the same anatomy with comparable
    intensities, so both the transformation and the similarity measure have an easy job. Different
    subjects imaged with different modalities combine two independent difficulties: the anatomies
    genuinely differ, and the intensities are not comparable. The relative difficulty of the two
    middle cases is a subjective judgement. Difficulty also depends on the anatomy --- brain is
    considerably easier than abdomen.
  \end{solution}
\answerspace{3.4cm}

  \part Each axis primarily determines one component of the registration algorithm. Which axis
  determines which, and why?

  \begin{solution}
    Answer: the \textbf{subject} axis determines the \textbf{transformation model}; the
    \textbf{modality} axis determines the \textbf{loss function}.

    Same subject means the same anatomy is present in both images, so a rigid or affine model is
    often enough --- what changed is position, not shape. Different subjects have genuinely
    different anatomy, with no exact correspondence and possibly pathology, so a deformable model
    is required.

    The modality determines whether intensities are comparable. Within a modality, SSD is
    appropriate. Across modalities it fails outright, and you need a statistical metric that is
    invariant to the intensity relationship --- normalised cross-correlation for a linear
    relationship, mutual information for an arbitrary one.
  \end{solution}
\answerspace{3.6cm}
\end{parts}

% =============================================================== Question 4 ===
\question
Please indicate whether the statement is true or false by circling the correct choice in each of
the following.

\begin{parts}
  \part \TFT\ \ The $3 \times 3$ mean filter is separable, $S_3 = \frac{1}{9}[1,1,1] * [1,1,1]^T$,
  which makes it cheaper to apply than a general $3 \times 3$ kernel.

  \part \TFF\ \ A Gaussian has finite support, so a $3 \times 3$ structural element represents a
  Gaussian filter exactly.

  \part \TFT\ \ Binomial filters are built by iterated convolution of $[1,1]$ and have better
  properties in the frequency domain than mean filtering.

  \part \TFF\ \ Increasing $\sigma$ of a Gaussian filter smooths more, but the size of structural
  element required is independent of $\sigma$.
\end{parts}

\begin{solution}
  (a) TRUE --- a separable 2D kernel can be applied as two 1D passes, which costs $O(k)$ per pixel
  instead of $O(k^2)$. The Gaussian and binomial kernels are separable as well.

  (b) FALSE --- the Gaussian has \emph{infinite} support. Any finite structural element is only an
  approximation of it, truncated at some distance from the centre.

  (c) TRUE --- $S_3 = \frac{1}{16}\left[\begin{smallmatrix}1&2&1\\2&4&2\\1&2&1\end{smallmatrix}
  \right]$, obtained from successive convolutions of $[1,1]$, and larger binomial kernels from
  successive convolutions of smaller ones. They are integer filters and behave better in the
  frequency domain than the box filter.

  (d) FALSE --- and (b) is the reason. Since the element only approximates the Gaussian by
  truncating it, a wider Gaussian needs a correspondingly larger element for the approximation to
  remain good. Using a $3 \times 3$ element with a large $\sigma$ silently truncates most of the
  kernel.
\end{solution}

% =============================================================== Question 5 ===
\question
A classification network takes an input image $x \in \mathbb{R}^{64 \times 64}$ with a single
channel and maps it through hidden layers of 256, 128 and 64 units to 16 output classes. Consider
only the \emph{first} link, from the input to the 256-unit layer, and ignore biases throughout.

\begin{parts}
  \part How many weights does this link have if it is fully connected?
  \begin{solution}
    Answer: $64 \times 64 \times 256 = 4\,096 \times 256 = \mathbf{1\,048\,576}$ weights.
  \end{solution}
\answerspace{1.8cm}

  \part How many weights does it have if it is instead a valid $5 \times 5$ convolution producing
  256 output channels from the single input channel?
  \begin{solution}
    Answer: $5 \times 5 \times 1 \times 256 = \mathbf{6\,400}$ weights.
  \end{solution}
\answerspace{1.8cm}

  \part What is the ratio, and what property of the convolutional layer explains it?
  \begin{solution}
    Answer: $1\,048\,576 / 6\,400 = \mathbf{163.84}$, so the fully connected link has roughly 164
    times more parameters.

    The explanation is \textbf{weight sharing}: the same $5 \times 5$ kernel is applied at every
    position, so the parameter count depends only on the kernel size and the channel counts --- and
    not at all on the image size. The fully connected count, by contrast, scales with the number of
    pixels. Doubling the image to $128 \times 128$ would quadruple the fully connected count to over
    four million while leaving the convolutional count at exactly 6\,400.

    Note also that the convolutional layer \emph{lifts} the data to 256 channels at every spatial
    position rather than compressing the whole image into 256 numbers, so it retains far more
    information despite using far fewer parameters.
  \end{solution}
\answerspace{3.4cm}
\end{parts}

% =============================================================== Question 6 ===
\question
In atlas-based segmentation the model is $p(I(x) \mid \theta) = \sum_n \mathcal{N}(I(x) \mid \mu_n,
\Sigma_n)\, p_{atlas}(c(x){=}n)$ with $\theta = \{\mu_n, \Sigma_n\}_{n=1}^N$. Please explain why the
intensity parameters $\mu_n$ and $\Sigma_n$ still have to be estimated for every new MR image, but
could in principle be fixed in advance for CT. Then state what else the atlas gives you compared
with plain EM. You do not need to write down the equations.

\begin{solution}
  Answer: \emph{Why MRI needs them fitted.} MRI has \textbf{no absolute intensity scale}. The number
  stored in a voxel depends on the scanner, the coil, the sequence parameters and the
  reconstruction, so ``the intensity of white matter'' is not a fixed quantity --- it differs
  between two scans of the same person on two machines. The Gaussian parameters therefore have to be
  re-estimated from each image.

  \emph{Why CT could fix them.} CT intensities are expressed in \textbf{Hounsfield units}, which are
  physically calibrated against water and air. A given tissue has a reproducible value across
  scanners, so $\mu_n$ and $\Sigma_n$ can be tabulated in advance and all the parameters fixed.

  \emph{What else the atlas buys.} Two things beyond the accuracy gain. It supplies
  \textbf{spatially varying prior information} --- at each voxel it says what is expected there,
  which is exactly what plain EM has no way of knowing. And it \textbf{removes the label
  ambiguity}: unsupervised EM returns clusters in an arbitrary order with no names attached, whereas
  with an atlas component $n$ \emph{is} grey matter. You know what you asked for, and there is no
  randomness left in which component is which.
\end{solution}
\answerspace{7.5cm}

% =============================================================== Question 7 ===
\question
Which of the following statements about interpolation is \textbf{inaccurate}?

\begin{choices}
  \choice Bilinear interpolation can equivalently be written as the parametric surface
  $I(x_1,x_2) = a + b x_1 + c x_2 + d x_1 x_2$, whose four unknowns are fixed by the intensities at
  the four surrounding pixels.
  \choice The trilinear extension adds the remaining monomials to give eight unknowns, fixed by the
  eight corner voxels of the cell containing the query point.
  \choice Bicubic interpolation has 16 unknowns, determined by the four corner intensities together
  with their first and mixed derivatives --- and since those derivatives must themselves be
  approximated from further pixels, bicubic effectively uses more than four samples.
  \choice Because bilinear interpolation is a weighted average of the surrounding values, applying
  it repeatedly during an iterative registration leaves the effective spatial resolution unchanged.
\end{choices}

\begin{solution}
  Answer: D.

  Averaging is exactly what costs you resolution. Bilinear interpolation smooths the image slightly
  every time it is applied, and during an optimisation loop it is applied thousands of times, so the
  blurring compounds. This is one reason the standard strategy is bilinear \emph{during} the
  optimisation, where speed matters, and bicubic for the \emph{final} resampling, where quality
  matters and the cost is paid once.
\end{solution}

% =============================================================== Question 8 ===
\question
Which of the following statements about the loss functions used for image restoration is
\textbf{incorrect}?

\begin{choices}
  \choice MAE, MSE, NMSE, PSNR and SSIM are all computed directly between the prediction and the
  ground-truth image.
  \choice SSIM compares many local patches of the two images and combines three components ---
  luminance, contrast and structure.
  \choice A perceptual loss passes both images through a previously trained network and compares
  intermediate features; that feature network must have been trained on the same task and the same
  modality as the restoration model, otherwise the comparison is meaningless.
  \choice An adversarial loss is a distributional distance defined over \emph{sets} of samples
  rather than between two paired images.
\end{choices}

\begin{solution}
  Answer: C.

  The feature network does not have to match the task or the modality. The lecture explicitly allows
  a well-established network trained on \emph{natural} images --- VGG is the standard choice --- or a
  network trained on CT images for some entirely different task. What is being borrowed is a
  general-purpose notion of what makes two images look different in terms of contextual content
  rather than pixel values, and that transfers surprisingly well across domains.

  The motivation for all of this is that MSE-type losses produce visibly over-smoothed output;
  combining a perceptual loss with an adversarial one gave the most realistic denoised low-dose CT
  in the referenced work.
\end{solution}

% =============================================================== Question 9 ===
\question
A semi-supervised auxiliary loss enforces consistency under spatial transformations. Two
transformations $\phi_1, \phi_2$ are drawn per image and the loss is
$\mathcal{L}\big(\phi_2 \circ \phi_1^{-1} \circ f(\phi_1 \circ x;\theta),\ f(\phi_2 \circ x;\theta)
\big)$. Which of the following statements is \textbf{incorrect}?

\begin{choices}
  \choice The loss requires no labels, so it can be evaluated on unlabelled images.
  \choice It encodes the principle that if $x \Leftrightarrow y$ then $\phi \circ x \Leftrightarrow
  \phi \circ y$, which is the same principle that underlies data augmentation.
  \choice The inverse $\phi_1^{-1}$ appears so that the two predictions are brought into a common
  spatial frame; without it they would sit in different frames and could not be compared pixel by
  pixel.
  \choice Because the loss is defined on unlabelled data, it replaces the supervised loss rather
  than being added to it.
\end{choices}

\begin{solution}
  Answer: D.

  It is an \emph{auxiliary} loss: the total cost is the supervised loss over the labelled examples
  \emph{plus} the consistency term, and the consistency term is summed over labelled and unlabelled
  examples alike. Dropping the supervised loss would remove the only thing that ties the network's
  output to the actual task --- a network that predicts a constant map everywhere is perfectly
  consistent under any transformation, and perfectly useless.
\end{solution}

% ============================================================== Question 10 ===
\question
When the network parameters are sampled from their posterior rather than optimised to a point
estimate, the sampled parameters can perform noticeably worse on the task than the MAP estimate.
What is the explanation given in the lecture?

\begin{choices}
  \choice In high dimensions the posterior is always unimodal, so every sample collapses onto the
  mean and no diversity is possible.
  \choice For an isotropic Gaussian in $d$ dimensions the expected distance of a sample from the
  mean grows like $\sqrt{d}$, so samples concentrate away from the mode --- and since the likelihood
  $p(D \mid \theta, M)$ \emph{is} the task loss, moving away from the mode means moving away from
  well-performing parameters.
  \choice High-dimensional posteriors cannot be sampled at all, so any apparent sample is an
  artefact of the random number generator.
  \choice In high dimensions the task loss becomes independent of the parameters, so performance is
  essentially random.
  \choice Sampling error decreases with dimension, so the effect disappears entirely for large
  networks.
\end{choices}

\begin{solution}
  Answer: B.

  This is the ``typical set'' phenomenon and it is thoroughly counter-intuitive: in high dimensions
  almost none of the probability mass sits near the mode, even though the mode is the single most
  probable point. Samples live in a thin shell at distance $\approx \sqrt{d}$ from the mean.

  The consequence for uncertainty estimation is uncomfortable. Because $\log p(D \mid \theta, M)$ is
  precisely the training objective, a parameter vector drawn honestly from the posterior can sit in
  a region of poor task performance --- so a faithful posterior sample and a well-performing network
  are not the same thing. Bounded sampling schemes have been proposed to mitigate this, but they do
  not fully solve it.
\end{solution}

% ============================================================== Question 11 ===
\question
A Vision Transformer is applied to images of size $224 \times 224$ with 3 channels. It uses a patch
size of $P = 16$, an embedding dimension of $D = 768$, and 12 attention heads. Recall that the image
is split into $N = HW/P^2$ patches, each flattened and mapped by a patch-embedding matrix
$E \in \mathbb{R}^{(P^2 C) \times D}$, and that a learnable class token is prepended to the sequence.
Circle True or False:

\begin{parts}
  \part \TFT\ \ The number of patch tokens is $N = 196$.

  \part \TFT\ \ The patch-embedding matrix $E$ is $768 \times 768$ and therefore contains
  $589\,824$ weights.

  \part \TFF\ \ After the class token is added, the sequence entering the encoder still has length
  196, because the class token takes the place of one patch.

  \part \TFF\ \ The attention matrix $QK^T$ of a single head has $N \times N$ entries, so its size
  grows linearly with the number of tokens.
\end{parts}

\begin{solution}
  (a) TRUE --- $N = HW/P^2 = (224 \times 224)/16^2 = 50\,176/256 = 196$, i.e.\ a $14 \times 14$ grid
  of patches.

  (b) TRUE --- the flattened patch has $P^2 C = 16^2 \times 3 = 256 \times 3 = 768$ entries, and it
  is mapped to $D = 768$, so $E$ is $768 \times 768 = 589\,824$. (That the two dimensions coincide is
  a property of this particular configuration, not a general rule.)

  (c) FALSE --- the class token is \emph{prepended}, not substituted: the sequence length becomes
  $196 + 1 = 197$. It carries no patch content of its own; it exists so that the final
  representation at that position can be read out and passed to the classification MLP.

  (d) FALSE --- an $N \times N$ matrix has $N^2$ entries, so the cost grows \textbf{quadratically}
  with sequence length. Here that is $197^2 = 38\,809$ attention scores per head, and 12 heads per
  layer. This quadratic scaling is the reason patching exists at all: attending over all
  $224 \times 224 = 50\,176$ pixels individually would need over $2.5$ billion scores per head.
\end{solution}

% ============================================================== Question 12 ===
\question
Which of the following statements on self-supervised learning are correct (multiple answers are
possible)?

\begin{choices}
  \choice Contrastive learning needs labelled examples in order to decide which pairs of images
  should be treated as positive and which as negative.
  \choice A pre-text task is a task that requires only images; the network trained on it is then
  fine-tuned with very few labelled examples on the task you actually care about.
  \choice In the \emph{local} contrastive loss the transformations must either preserve the patch
  correspondences or be undone before the corresponding patches are compared.
  \choice Foundation models in this setting are obtained by supervised training at very large scale
  on carefully curated labelled data sets.
  \choice The InfoNCE loss maximises the similarity between two views of the same image through its
  numerator, and minimises the similarity to the samples of the negative set through its
  denominator.
\end{choices}

\begin{solution}
  Answer: (B), (C) and (E) are correct.

  (A) is wrong, and it is wrong about the whole point of the method. The labels in contrastive
  learning are \emph{surrogate} labels manufactured from the data itself: two random transformations
  of the same image are declared ``close'', everything else ``far''. No human annotation is involved
  --- which is precisely why the approach is useful when labels are scarce. Note also that the
  transformations are not restricted to geometric ones; intensity changes are used deliberately, so
  that the representation becomes invariant to them.

  (D) inverts the defining property. Foundation models here are built by \textbf{self-supervised}
  training at very large scale --- masked auto-encoding and variations of contrastive learning ---
  on transformer architectures. Requiring curated labels at that scale is exactly what these methods
  exist to avoid.
\end{solution}

\end{questions}

\end{document}
"
%% Compile the blank version (no answers):  pdflatex Sample_Exam_6.tex
%%
\ifx\withanswers\undefined
  \documentclass[11pt]{exam}
\else
  \documentclass[11pt,answers]{exam}
\fi

\usepackage[margin=2.4cm]{geometry}
\usepackage{amsmath,amssymb}
\usepackage{array}
\usepackage[T1]{fontenc}

% ---------------------------------------------------------------- page style
\pagestyle{foot}
\firstpageheader{}{}{}
\firstpagefooter{}{}{}
\runningheader{}{}{}
\runningfooter{}{}{Page \thepage}

% -------------------------------------------------------------- multiplechoice
% Reproduce the "(A) (B) (C) ..." labelling of the original exam.
\renewcommand{\choicelabel}{$\bigcirc$~(\Alph{choice})}

% ---------------------------------------------------------- true/false macros
% \TFT -> the statement is TRUE, \TFF -> the statement is FALSE.
% In the answer version the correct word is set in capitals and bold.
\newcommand{\TFT}{\ifprintanswers\textbf{TRUE}\hspace{1.1em}False\else True\hspace{1.1em}False\fi}
\newcommand{\TFF}{True\hspace{1.1em}\ifprintanswers\textbf{FALSE}\else False\fi}

% ------------------------------------------------------------- writing space
% Reserves room to write in the blank version; collapses in the answer version.
\newcommand{\answerspace}[1]{\ifprintanswers\else\vspace{#1}\fi}

\begin{document}

% ------------------------------------------------------------------ title page
\begin{center}
  {\LARGE\bfseries Sample Exam}\\[2mm]
  {\Large Medical Image Analysis}\\[6mm]
\end{center}

\vspace{4mm}
\noindent Student name:\hrulefill

\vspace{6mm}
\noindent Immatriculation number:\hrulefill

\vspace{10mm}

\begin{center}
  {\large\bfseries Exam Questions}
\end{center}

\noindent\textbf{Note:} All the questions have equal weight.

\vspace{4mm}

\begin{questions}

% =============================================================== Question 1 ===
\question
In axial 2D acquisitions the slices lie in the transverse plane and are stacked along the
head-to-toe direction. Voxel sizes are given as in-plane (two values) and through-plane (the slice
thickness), and the field-of-view is image size $\times$ voxel size.

\vspace{1mm}
\noindent An axial acquisition uses in-plane voxel size $0.7 \times 0.7$~mm$^2$ with 320 pixels in
both directions, and through-plane voxel size $4.0$~mm with 44 slices.

\begin{parts}

  \part What is the field-of-view in the head-to-toe, left-to-right and anterior-to-posterior
  directions, respectively?

  \begin{solution}
    Answer: $176 \times 224 \times 224$~mm$^3$.

    Head-to-toe is the through-plane direction of an axial acquisition: $44 \times 4.0 = 176$~mm.
    Both in-plane directions give $320 \times 0.7 = 224$~mm.
  \end{solution}
\answerspace{2.2cm}

  \part The anatomy that has to be covered extends 210~mm in the head-to-toe direction. Does this
  protocol cover it? If not, how many slices are needed at the same slice thickness, and by what
  factor does the acquisition time change if it is proportional to the number of slices?

  \begin{solution}
    Answer: \textbf{No} --- the protocol covers only 176~mm of the required 210~mm.
    $\lceil 210 / 4.0 \rceil = \lceil 52.5 \rceil = \mathbf{53}$ slices are needed (53 slices cover
    212~mm). The acquisition time grows by a factor $53/44 \approx \mathbf{1.20}$, i.e.\ about
    20\% longer.

    Note that you must round \emph{up}: 52 slices would cover only 208~mm and miss the anatomy.
  \end{solution}
\answerspace{3.2cm}

  \part Compute the anisotropy ratio of the voxels. This volume is then fed, without resampling,
  to a 3D convolutional network that was trained on $1.0$~mm isotropic data. What goes wrong?

  \begin{solution}
    Answer: the anisotropy ratio is $4.0 / 0.7 \approx \mathbf{5.7}$ --- the voxel is almost six
    times longer through-plane than in-plane.

    A $3 \times 3 \times 3$ kernel therefore covers $2.1 \times 2.1 \times 12.0$~mm$^3$ of real
    tissue instead of the $3 \times 3 \times 3$~mm$^3$ it saw during training. The receptive field
    of every neuron is wildly asymmetric in physical space, so the spatial features the network
    learned simply do not correspond to the same physical structures any more. This is a domain
    shift the model was never trained for, and the fix is to resample to the training voxel size
    first.

    Note that resampling does not create the missing information: upsampling the through-plane
    direction from 4.0~mm to 1.0~mm interpolates three samples out of every one that was really
    measured. The effective resolution in that direction is still 4~mm.
  \end{solution}
\answerspace{4.5cm}

\end{parts}

% =============================================================== Question 2 ===
\question
In atlas-based segmentation the atlas is registered to the image and then used inside the mixture
model. Which of the two posteriors below is the one used by atlas-based segmentation, as opposed to
plain EM with a Gaussian mixture model? $\pi_n$ denotes a mixture weight estimated from the image
itself. Please circle the correct choice.

\vspace{2mm}
\begin{center}
  $\displaystyle p(c(x){=}n \mid I(x)) \propto \mathcal{N}(I(x) \mid \mu_n, \Sigma_n)\ \pi_n$
  \hspace{1.2cm}
  $\displaystyle p(c(x){=}n \mid I(x)) \propto \mathcal{N}(I(x) \mid \mu_n, \Sigma_n)\
  p_{atlas}(c(x){=}n)$
\end{center}
\vspace{2mm}

\begin{solution}
  Answer: Right one.

  The atlas prior $p_{atlas}(c(x){=}n)$ \emph{replaces} the free mixture weight $\pi_n$. Two things
  change as a result. The prior is now spatially varying --- it says something different at every
  voxel, which is exactly the information plain EM lacks --- and it no longer has to be estimated,
  so the class probabilities drop out of the optimisation entirely.
\end{solution}
\answerspace{1.5cm}

% =============================================================== Question 3 ===
\question
Registration problems are classified along two axes: whether the two images come from the
\emph{same or different subjects}, and whether they were acquired with the \emph{same or different
modalities}.

\begin{parts}
  \part Which of the four combinations is the easiest and which is the hardest, and why?

  \begin{solution}
    Answer: easiest is \textbf{intra-subject, intra-modality}; hardest is \textbf{inter-subject,
    inter-modality}.

    The same subject imaged twice with the same sequence shows the same anatomy with comparable
    intensities, so both the transformation and the similarity measure have an easy job. Different
    subjects imaged with different modalities combine two independent difficulties: the anatomies
    genuinely differ, and the intensities are not comparable. The relative difficulty of the two
    middle cases is a subjective judgement. Difficulty also depends on the anatomy --- brain is
    considerably easier than abdomen.
  \end{solution}
\answerspace{3.4cm}

  \part Each axis primarily determines one component of the registration algorithm. Which axis
  determines which, and why?

  \begin{solution}
    Answer: the \textbf{subject} axis determines the \textbf{transformation model}; the
    \textbf{modality} axis determines the \textbf{loss function}.

    Same subject means the same anatomy is present in both images, so a rigid or affine model is
    often enough --- what changed is position, not shape. Different subjects have genuinely
    different anatomy, with no exact correspondence and possibly pathology, so a deformable model
    is required.

    The modality determines whether intensities are comparable. Within a modality, SSD is
    appropriate. Across modalities it fails outright, and you need a statistical metric that is
    invariant to the intensity relationship --- normalised cross-correlation for a linear
    relationship, mutual information for an arbitrary one.
  \end{solution}
\answerspace{3.6cm}
\end{parts}

% =============================================================== Question 4 ===
\question
Please indicate whether the statement is true or false by circling the correct choice in each of
the following.

\begin{parts}
  \part \TFT\ \ The $3 \times 3$ mean filter is separable, $S_3 = \frac{1}{9}[1,1,1] * [1,1,1]^T$,
  which makes it cheaper to apply than a general $3 \times 3$ kernel.

  \part \TFF\ \ A Gaussian has finite support, so a $3 \times 3$ structural element represents a
  Gaussian filter exactly.

  \part \TFT\ \ Binomial filters are built by iterated convolution of $[1,1]$ and have better
  properties in the frequency domain than mean filtering.

  \part \TFF\ \ Increasing $\sigma$ of a Gaussian filter smooths more, but the size of structural
  element required is independent of $\sigma$.
\end{parts}

\begin{solution}
  (a) TRUE --- a separable 2D kernel can be applied as two 1D passes, which costs $O(k)$ per pixel
  instead of $O(k^2)$. The Gaussian and binomial kernels are separable as well.

  (b) FALSE --- the Gaussian has \emph{infinite} support. Any finite structural element is only an
  approximation of it, truncated at some distance from the centre.

  (c) TRUE --- $S_3 = \frac{1}{16}\left[\begin{smallmatrix}1&2&1\\2&4&2\\1&2&1\end{smallmatrix}
  \right]$, obtained from successive convolutions of $[1,1]$, and larger binomial kernels from
  successive convolutions of smaller ones. They are integer filters and behave better in the
  frequency domain than the box filter.

  (d) FALSE --- and (b) is the reason. Since the element only approximates the Gaussian by
  truncating it, a wider Gaussian needs a correspondingly larger element for the approximation to
  remain good. Using a $3 \times 3$ element with a large $\sigma$ silently truncates most of the
  kernel.
\end{solution}

% =============================================================== Question 5 ===
\question
A classification network takes an input image $x \in \mathbb{R}^{64 \times 64}$ with a single
channel and maps it through hidden layers of 256, 128 and 64 units to 16 output classes. Consider
only the \emph{first} link, from the input to the 256-unit layer, and ignore biases throughout.

\begin{parts}
  \part How many weights does this link have if it is fully connected?
  \begin{solution}
    Answer: $64 \times 64 \times 256 = 4\,096 \times 256 = \mathbf{1\,048\,576}$ weights.
  \end{solution}
\answerspace{1.8cm}

  \part How many weights does it have if it is instead a valid $5 \times 5$ convolution producing
  256 output channels from the single input channel?
  \begin{solution}
    Answer: $5 \times 5 \times 1 \times 256 = \mathbf{6\,400}$ weights.
  \end{solution}
\answerspace{1.8cm}

  \part What is the ratio, and what property of the convolutional layer explains it?
  \begin{solution}
    Answer: $1\,048\,576 / 6\,400 = \mathbf{163.84}$, so the fully connected link has roughly 164
    times more parameters.

    The explanation is \textbf{weight sharing}: the same $5 \times 5$ kernel is applied at every
    position, so the parameter count depends only on the kernel size and the channel counts --- and
    not at all on the image size. The fully connected count, by contrast, scales with the number of
    pixels. Doubling the image to $128 \times 128$ would quadruple the fully connected count to over
    four million while leaving the convolutional count at exactly 6\,400.

    Note also that the convolutional layer \emph{lifts} the data to 256 channels at every spatial
    position rather than compressing the whole image into 256 numbers, so it retains far more
    information despite using far fewer parameters.
  \end{solution}
\answerspace{3.4cm}
\end{parts}

% =============================================================== Question 6 ===
\question
In atlas-based segmentation the model is $p(I(x) \mid \theta) = \sum_n \mathcal{N}(I(x) \mid \mu_n,
\Sigma_n)\, p_{atlas}(c(x){=}n)$ with $\theta = \{\mu_n, \Sigma_n\}_{n=1}^N$. Please explain why the
intensity parameters $\mu_n$ and $\Sigma_n$ still have to be estimated for every new MR image, but
could in principle be fixed in advance for CT. Then state what else the atlas gives you compared
with plain EM. You do not need to write down the equations.

\begin{solution}
  Answer: \emph{Why MRI needs them fitted.} MRI has \textbf{no absolute intensity scale}. The number
  stored in a voxel depends on the scanner, the coil, the sequence parameters and the
  reconstruction, so ``the intensity of white matter'' is not a fixed quantity --- it differs
  between two scans of the same person on two machines. The Gaussian parameters therefore have to be
  re-estimated from each image.

  \emph{Why CT could fix them.} CT intensities are expressed in \textbf{Hounsfield units}, which are
  physically calibrated against water and air. A given tissue has a reproducible value across
  scanners, so $\mu_n$ and $\Sigma_n$ can be tabulated in advance and all the parameters fixed.

  \emph{What else the atlas buys.} Two things beyond the accuracy gain. It supplies
  \textbf{spatially varying prior information} --- at each voxel it says what is expected there,
  which is exactly what plain EM has no way of knowing. And it \textbf{removes the label
  ambiguity}: unsupervised EM returns clusters in an arbitrary order with no names attached, whereas
  with an atlas component $n$ \emph{is} grey matter. You know what you asked for, and there is no
  randomness left in which component is which.
\end{solution}
\answerspace{7.5cm}

% =============================================================== Question 7 ===
\question
Which of the following statements about interpolation is \textbf{inaccurate}?

\begin{choices}
  \choice Bilinear interpolation can equivalently be written as the parametric surface
  $I(x_1,x_2) = a + b x_1 + c x_2 + d x_1 x_2$, whose four unknowns are fixed by the intensities at
  the four surrounding pixels.
  \choice The trilinear extension adds the remaining monomials to give eight unknowns, fixed by the
  eight corner voxels of the cell containing the query point.
  \choice Bicubic interpolation has 16 unknowns, determined by the four corner intensities together
  with their first and mixed derivatives --- and since those derivatives must themselves be
  approximated from further pixels, bicubic effectively uses more than four samples.
  \choice Because bilinear interpolation is a weighted average of the surrounding values, applying
  it repeatedly during an iterative registration leaves the effective spatial resolution unchanged.
\end{choices}

\begin{solution}
  Answer: D.

  Averaging is exactly what costs you resolution. Bilinear interpolation smooths the image slightly
  every time it is applied, and during an optimisation loop it is applied thousands of times, so the
  blurring compounds. This is one reason the standard strategy is bilinear \emph{during} the
  optimisation, where speed matters, and bicubic for the \emph{final} resampling, where quality
  matters and the cost is paid once.
\end{solution}

% =============================================================== Question 8 ===
\question
Which of the following statements about the loss functions used for image restoration is
\textbf{incorrect}?

\begin{choices}
  \choice MAE, MSE, NMSE, PSNR and SSIM are all computed directly between the prediction and the
  ground-truth image.
  \choice SSIM compares many local patches of the two images and combines three components ---
  luminance, contrast and structure.
  \choice A perceptual loss passes both images through a previously trained network and compares
  intermediate features; that feature network must have been trained on the same task and the same
  modality as the restoration model, otherwise the comparison is meaningless.
  \choice An adversarial loss is a distributional distance defined over \emph{sets} of samples
  rather than between two paired images.
\end{choices}

\begin{solution}
  Answer: C.

  The feature network does not have to match the task or the modality. The lecture explicitly allows
  a well-established network trained on \emph{natural} images --- VGG is the standard choice --- or a
  network trained on CT images for some entirely different task. What is being borrowed is a
  general-purpose notion of what makes two images look different in terms of contextual content
  rather than pixel values, and that transfers surprisingly well across domains.

  The motivation for all of this is that MSE-type losses produce visibly over-smoothed output;
  combining a perceptual loss with an adversarial one gave the most realistic denoised low-dose CT
  in the referenced work.
\end{solution}

% =============================================================== Question 9 ===
\question
A semi-supervised auxiliary loss enforces consistency under spatial transformations. Two
transformations $\phi_1, \phi_2$ are drawn per image and the loss is
$\mathcal{L}\big(\phi_2 \circ \phi_1^{-1} \circ f(\phi_1 \circ x;\theta),\ f(\phi_2 \circ x;\theta)
\big)$. Which of the following statements is \textbf{incorrect}?

\begin{choices}
  \choice The loss requires no labels, so it can be evaluated on unlabelled images.
  \choice It encodes the principle that if $x \Leftrightarrow y$ then $\phi \circ x \Leftrightarrow
  \phi \circ y$, which is the same principle that underlies data augmentation.
  \choice The inverse $\phi_1^{-1}$ appears so that the two predictions are brought into a common
  spatial frame; without it they would sit in different frames and could not be compared pixel by
  pixel.
  \choice Because the loss is defined on unlabelled data, it replaces the supervised loss rather
  than being added to it.
\end{choices}

\begin{solution}
  Answer: D.

  It is an \emph{auxiliary} loss: the total cost is the supervised loss over the labelled examples
  \emph{plus} the consistency term, and the consistency term is summed over labelled and unlabelled
  examples alike. Dropping the supervised loss would remove the only thing that ties the network's
  output to the actual task --- a network that predicts a constant map everywhere is perfectly
  consistent under any transformation, and perfectly useless.
\end{solution}

% ============================================================== Question 10 ===
\question
When the network parameters are sampled from their posterior rather than optimised to a point
estimate, the sampled parameters can perform noticeably worse on the task than the MAP estimate.
What is the explanation given in the lecture?

\begin{choices}
  \choice In high dimensions the posterior is always unimodal, so every sample collapses onto the
  mean and no diversity is possible.
  \choice For an isotropic Gaussian in $d$ dimensions the expected distance of a sample from the
  mean grows like $\sqrt{d}$, so samples concentrate away from the mode --- and since the likelihood
  $p(D \mid \theta, M)$ \emph{is} the task loss, moving away from the mode means moving away from
  well-performing parameters.
  \choice High-dimensional posteriors cannot be sampled at all, so any apparent sample is an
  artefact of the random number generator.
  \choice In high dimensions the task loss becomes independent of the parameters, so performance is
  essentially random.
  \choice Sampling error decreases with dimension, so the effect disappears entirely for large
  networks.
\end{choices}

\begin{solution}
  Answer: B.

  This is the ``typical set'' phenomenon and it is thoroughly counter-intuitive: in high dimensions
  almost none of the probability mass sits near the mode, even though the mode is the single most
  probable point. Samples live in a thin shell at distance $\approx \sqrt{d}$ from the mean.

  The consequence for uncertainty estimation is uncomfortable. Because $\log p(D \mid \theta, M)$ is
  precisely the training objective, a parameter vector drawn honestly from the posterior can sit in
  a region of poor task performance --- so a faithful posterior sample and a well-performing network
  are not the same thing. Bounded sampling schemes have been proposed to mitigate this, but they do
  not fully solve it.
\end{solution}

% ============================================================== Question 11 ===
\question
A Vision Transformer is applied to images of size $224 \times 224$ with 3 channels. It uses a patch
size of $P = 16$, an embedding dimension of $D = 768$, and 12 attention heads. Recall that the image
is split into $N = HW/P^2$ patches, each flattened and mapped by a patch-embedding matrix
$E \in \mathbb{R}^{(P^2 C) \times D}$, and that a learnable class token is prepended to the sequence.
Circle True or False:

\begin{parts}
  \part \TFT\ \ The number of patch tokens is $N = 196$.

  \part \TFT\ \ The patch-embedding matrix $E$ is $768 \times 768$ and therefore contains
  $589\,824$ weights.

  \part \TFF\ \ After the class token is added, the sequence entering the encoder still has length
  196, because the class token takes the place of one patch.

  \part \TFF\ \ The attention matrix $QK^T$ of a single head has $N \times N$ entries, so its size
  grows linearly with the number of tokens.
\end{parts}

\begin{solution}
  (a) TRUE --- $N = HW/P^2 = (224 \times 224)/16^2 = 50\,176/256 = 196$, i.e.\ a $14 \times 14$ grid
  of patches.

  (b) TRUE --- the flattened patch has $P^2 C = 16^2 \times 3 = 256 \times 3 = 768$ entries, and it
  is mapped to $D = 768$, so $E$ is $768 \times 768 = 589\,824$. (That the two dimensions coincide is
  a property of this particular configuration, not a general rule.)

  (c) FALSE --- the class token is \emph{prepended}, not substituted: the sequence length becomes
  $196 + 1 = 197$. It carries no patch content of its own; it exists so that the final
  representation at that position can be read out and passed to the classification MLP.

  (d) FALSE --- an $N \times N$ matrix has $N^2$ entries, so the cost grows \textbf{quadratically}
  with sequence length. Here that is $197^2 = 38\,809$ attention scores per head, and 12 heads per
  layer. This quadratic scaling is the reason patching exists at all: attending over all
  $224 \times 224 = 50\,176$ pixels individually would need over $2.5$ billion scores per head.
\end{solution}

% ============================================================== Question 12 ===
\question
Which of the following statements on self-supervised learning are correct (multiple answers are
possible)?

\begin{choices}
  \choice Contrastive learning needs labelled examples in order to decide which pairs of images
  should be treated as positive and which as negative.
  \choice A pre-text task is a task that requires only images; the network trained on it is then
  fine-tuned with very few labelled examples on the task you actually care about.
  \choice In the \emph{local} contrastive loss the transformations must either preserve the patch
  correspondences or be undone before the corresponding patches are compared.
  \choice Foundation models in this setting are obtained by supervised training at very large scale
  on carefully curated labelled data sets.
  \choice The InfoNCE loss maximises the similarity between two views of the same image through its
  numerator, and minimises the similarity to the samples of the negative set through its
  denominator.
\end{choices}

\begin{solution}
  Answer: (B), (C) and (E) are correct.

  (A) is wrong, and it is wrong about the whole point of the method. The labels in contrastive
  learning are \emph{surrogate} labels manufactured from the data itself: two random transformations
  of the same image are declared ``close'', everything else ``far''. No human annotation is involved
  --- which is precisely why the approach is useful when labels are scarce. Note also that the
  transformations are not restricted to geometric ones; intensity changes are used deliberately, so
  that the representation becomes invariant to them.

  (D) inverts the defining property. Foundation models here are built by \textbf{self-supervised}
  training at very large scale --- masked auto-encoding and variations of contrastive learning ---
  on transformer architectures. Requiring curated labels at that scale is exactly what these methods
  exist to avoid.
\end{solution}

\end{questions}

\end{document}
"
%% Compile the blank version (no answers):  pdflatex Sample_Exam_6.tex
%%
\ifx\withanswers\undefined
  \documentclass[11pt]{exam}
\else
  \documentclass[11pt,answers]{exam}
\fi

\usepackage[margin=2.4cm]{geometry}
\usepackage{amsmath,amssymb}
\usepackage{array}
\usepackage[T1]{fontenc}

% ---------------------------------------------------------------- page style
\pagestyle{foot}
\firstpageheader{}{}{}
\firstpagefooter{}{}{}
\runningheader{}{}{}
\runningfooter{}{}{Page \thepage}

% -------------------------------------------------------------- multiplechoice
% Reproduce the "(A) (B) (C) ..." labelling of the original exam.
\renewcommand{\choicelabel}{$\bigcirc$~(\Alph{choice})}

% ---------------------------------------------------------- true/false macros
% \TFT -> the statement is TRUE, \TFF -> the statement is FALSE.
% In the answer version the correct word is set in capitals and bold.
\newcommand{\TFT}{\ifprintanswers\textbf{TRUE}\hspace{1.1em}False\else True\hspace{1.1em}False\fi}
\newcommand{\TFF}{True\hspace{1.1em}\ifprintanswers\textbf{FALSE}\else False\fi}

% ------------------------------------------------------------- writing space
% Reserves room to write in the blank version; collapses in the answer version.
\newcommand{\answerspace}[1]{\ifprintanswers\else\vspace{#1}\fi}

\begin{document}

% ------------------------------------------------------------------ title page
\begin{center}
  {\LARGE\bfseries Sample Exam}\\[2mm]
  {\Large Medical Image Analysis}\\[6mm]
\end{center}

\vspace{4mm}
\noindent Student name:\hrulefill

\vspace{6mm}
\noindent Immatriculation number:\hrulefill

\vspace{10mm}

\begin{center}
  {\large\bfseries Exam Questions}
\end{center}

\noindent\textbf{Note:} All the questions have equal weight.

\vspace{4mm}

\begin{questions}

% =============================================================== Question 1 ===
\question
In axial 2D acquisitions the slices lie in the transverse plane and are stacked along the
head-to-toe direction. Voxel sizes are given as in-plane (two values) and through-plane (the slice
thickness), and the field-of-view is image size $\times$ voxel size.

\vspace{1mm}
\noindent An axial acquisition uses in-plane voxel size $0.7 \times 0.7$~mm$^2$ with 320 pixels in
both directions, and through-plane voxel size $4.0$~mm with 44 slices.

\begin{parts}

  \part What is the field-of-view in the head-to-toe, left-to-right and anterior-to-posterior
  directions, respectively?

  \begin{solution}
    Answer: $176 \times 224 \times 224$~mm$^3$.

    Head-to-toe is the through-plane direction of an axial acquisition: $44 \times 4.0 = 176$~mm.
    Both in-plane directions give $320 \times 0.7 = 224$~mm.
  \end{solution}
\answerspace{2.2cm}

  \part The anatomy that has to be covered extends 210~mm in the head-to-toe direction. Does this
  protocol cover it? If not, how many slices are needed at the same slice thickness, and by what
  factor does the acquisition time change if it is proportional to the number of slices?

  \begin{solution}
    Answer: \textbf{No} --- the protocol covers only 176~mm of the required 210~mm.
    $\lceil 210 / 4.0 \rceil = \lceil 52.5 \rceil = \mathbf{53}$ slices are needed (53 slices cover
    212~mm). The acquisition time grows by a factor $53/44 \approx \mathbf{1.20}$, i.e.\ about
    20\% longer.

    Note that you must round \emph{up}: 52 slices would cover only 208~mm and miss the anatomy.
  \end{solution}
\answerspace{3.2cm}

  \part Compute the anisotropy ratio of the voxels. This volume is then fed, without resampling,
  to a 3D convolutional network that was trained on $1.0$~mm isotropic data. What goes wrong?

  \begin{solution}
    Answer: the anisotropy ratio is $4.0 / 0.7 \approx \mathbf{5.7}$ --- the voxel is almost six
    times longer through-plane than in-plane.

    A $3 \times 3 \times 3$ kernel therefore covers $2.1 \times 2.1 \times 12.0$~mm$^3$ of real
    tissue instead of the $3 \times 3 \times 3$~mm$^3$ it saw during training. The receptive field
    of every neuron is wildly asymmetric in physical space, so the spatial features the network
    learned simply do not correspond to the same physical structures any more. This is a domain
    shift the model was never trained for, and the fix is to resample to the training voxel size
    first.

    Note that resampling does not create the missing information: upsampling the through-plane
    direction from 4.0~mm to 1.0~mm interpolates three samples out of every one that was really
    measured. The effective resolution in that direction is still 4~mm.
  \end{solution}
\answerspace{4.5cm}

\end{parts}

% =============================================================== Question 2 ===
\question
In atlas-based segmentation the atlas is registered to the image and then used inside the mixture
model. Which of the two posteriors below is the one used by atlas-based segmentation, as opposed to
plain EM with a Gaussian mixture model? $\pi_n$ denotes a mixture weight estimated from the image
itself. Please circle the correct choice.

\vspace{2mm}
\begin{center}
  $\displaystyle p(c(x){=}n \mid I(x)) \propto \mathcal{N}(I(x) \mid \mu_n, \Sigma_n)\ \pi_n$
  \hspace{1.2cm}
  $\displaystyle p(c(x){=}n \mid I(x)) \propto \mathcal{N}(I(x) \mid \mu_n, \Sigma_n)\
  p_{atlas}(c(x){=}n)$
\end{center}
\vspace{2mm}

\begin{solution}
  Answer: Right one.

  The atlas prior $p_{atlas}(c(x){=}n)$ \emph{replaces} the free mixture weight $\pi_n$. Two things
  change as a result. The prior is now spatially varying --- it says something different at every
  voxel, which is exactly the information plain EM lacks --- and it no longer has to be estimated,
  so the class probabilities drop out of the optimisation entirely.
\end{solution}
\answerspace{1.5cm}

% =============================================================== Question 3 ===
\question
Registration problems are classified along two axes: whether the two images come from the
\emph{same or different subjects}, and whether they were acquired with the \emph{same or different
modalities}.

\begin{parts}
  \part Which of the four combinations is the easiest and which is the hardest, and why?

  \begin{solution}
    Answer: easiest is \textbf{intra-subject, intra-modality}; hardest is \textbf{inter-subject,
    inter-modality}.

    The same subject imaged twice with the same sequence shows the same anatomy with comparable
    intensities, so both the transformation and the similarity measure have an easy job. Different
    subjects imaged with different modalities combine two independent difficulties: the anatomies
    genuinely differ, and the intensities are not comparable. The relative difficulty of the two
    middle cases is a subjective judgement. Difficulty also depends on the anatomy --- brain is
    considerably easier than abdomen.
  \end{solution}
\answerspace{3.4cm}

  \part Each axis primarily determines one component of the registration algorithm. Which axis
  determines which, and why?

  \begin{solution}
    Answer: the \textbf{subject} axis determines the \textbf{transformation model}; the
    \textbf{modality} axis determines the \textbf{loss function}.

    Same subject means the same anatomy is present in both images, so a rigid or affine model is
    often enough --- what changed is position, not shape. Different subjects have genuinely
    different anatomy, with no exact correspondence and possibly pathology, so a deformable model
    is required.

    The modality determines whether intensities are comparable. Within a modality, SSD is
    appropriate. Across modalities it fails outright, and you need a statistical metric that is
    invariant to the intensity relationship --- normalised cross-correlation for a linear
    relationship, mutual information for an arbitrary one.
  \end{solution}
\answerspace{3.6cm}
\end{parts}

% =============================================================== Question 4 ===
\question
Please indicate whether the statement is true or false by circling the correct choice in each of
the following.

\begin{parts}
  \part \TFT\ \ The $3 \times 3$ mean filter is separable, $S_3 = \frac{1}{9}[1,1,1] * [1,1,1]^T$,
  which makes it cheaper to apply than a general $3 \times 3$ kernel.

  \part \TFF\ \ A Gaussian has finite support, so a $3 \times 3$ structural element represents a
  Gaussian filter exactly.

  \part \TFT\ \ Binomial filters are built by iterated convolution of $[1,1]$ and have better
  properties in the frequency domain than mean filtering.

  \part \TFF\ \ Increasing $\sigma$ of a Gaussian filter smooths more, but the size of structural
  element required is independent of $\sigma$.
\end{parts}

\begin{solution}
  (a) TRUE --- a separable 2D kernel can be applied as two 1D passes, which costs $O(k)$ per pixel
  instead of $O(k^2)$. The Gaussian and binomial kernels are separable as well.

  (b) FALSE --- the Gaussian has \emph{infinite} support. Any finite structural element is only an
  approximation of it, truncated at some distance from the centre.

  (c) TRUE --- $S_3 = \frac{1}{16}\left[\begin{smallmatrix}1&2&1\\2&4&2\\1&2&1\end{smallmatrix}
  \right]$, obtained from successive convolutions of $[1,1]$, and larger binomial kernels from
  successive convolutions of smaller ones. They are integer filters and behave better in the
  frequency domain than the box filter.

  (d) FALSE --- and (b) is the reason. Since the element only approximates the Gaussian by
  truncating it, a wider Gaussian needs a correspondingly larger element for the approximation to
  remain good. Using a $3 \times 3$ element with a large $\sigma$ silently truncates most of the
  kernel.
\end{solution}

% =============================================================== Question 5 ===
\question
A classification network takes an input image $x \in \mathbb{R}^{64 \times 64}$ with a single
channel and maps it through hidden layers of 256, 128 and 64 units to 16 output classes. Consider
only the \emph{first} link, from the input to the 256-unit layer, and ignore biases throughout.

\begin{parts}
  \part How many weights does this link have if it is fully connected?
  \begin{solution}
    Answer: $64 \times 64 \times 256 = 4\,096 \times 256 = \mathbf{1\,048\,576}$ weights.
  \end{solution}
\answerspace{1.8cm}

  \part How many weights does it have if it is instead a valid $5 \times 5$ convolution producing
  256 output channels from the single input channel?
  \begin{solution}
    Answer: $5 \times 5 \times 1 \times 256 = \mathbf{6\,400}$ weights.
  \end{solution}
\answerspace{1.8cm}

  \part What is the ratio, and what property of the convolutional layer explains it?
  \begin{solution}
    Answer: $1\,048\,576 / 6\,400 = \mathbf{163.84}$, so the fully connected link has roughly 164
    times more parameters.

    The explanation is \textbf{weight sharing}: the same $5 \times 5$ kernel is applied at every
    position, so the parameter count depends only on the kernel size and the channel counts --- and
    not at all on the image size. The fully connected count, by contrast, scales with the number of
    pixels. Doubling the image to $128 \times 128$ would quadruple the fully connected count to over
    four million while leaving the convolutional count at exactly 6\,400.

    Note also that the convolutional layer \emph{lifts} the data to 256 channels at every spatial
    position rather than compressing the whole image into 256 numbers, so it retains far more
    information despite using far fewer parameters.
  \end{solution}
\answerspace{3.4cm}
\end{parts}

% =============================================================== Question 6 ===
\question
In atlas-based segmentation the model is $p(I(x) \mid \theta) = \sum_n \mathcal{N}(I(x) \mid \mu_n,
\Sigma_n)\, p_{atlas}(c(x){=}n)$ with $\theta = \{\mu_n, \Sigma_n\}_{n=1}^N$. Please explain why the
intensity parameters $\mu_n$ and $\Sigma_n$ still have to be estimated for every new MR image, but
could in principle be fixed in advance for CT. Then state what else the atlas gives you compared
with plain EM. You do not need to write down the equations.

\begin{solution}
  Answer: \emph{Why MRI needs them fitted.} MRI has \textbf{no absolute intensity scale}. The number
  stored in a voxel depends on the scanner, the coil, the sequence parameters and the
  reconstruction, so ``the intensity of white matter'' is not a fixed quantity --- it differs
  between two scans of the same person on two machines. The Gaussian parameters therefore have to be
  re-estimated from each image.

  \emph{Why CT could fix them.} CT intensities are expressed in \textbf{Hounsfield units}, which are
  physically calibrated against water and air. A given tissue has a reproducible value across
  scanners, so $\mu_n$ and $\Sigma_n$ can be tabulated in advance and all the parameters fixed.

  \emph{What else the atlas buys.} Two things beyond the accuracy gain. It supplies
  \textbf{spatially varying prior information} --- at each voxel it says what is expected there,
  which is exactly what plain EM has no way of knowing. And it \textbf{removes the label
  ambiguity}: unsupervised EM returns clusters in an arbitrary order with no names attached, whereas
  with an atlas component $n$ \emph{is} grey matter. You know what you asked for, and there is no
  randomness left in which component is which.
\end{solution}
\answerspace{7.5cm}

% =============================================================== Question 7 ===
\question
Which of the following statements about interpolation is \textbf{inaccurate}?

\begin{choices}
  \choice Bilinear interpolation can equivalently be written as the parametric surface
  $I(x_1,x_2) = a + b x_1 + c x_2 + d x_1 x_2$, whose four unknowns are fixed by the intensities at
  the four surrounding pixels.
  \choice The trilinear extension adds the remaining monomials to give eight unknowns, fixed by the
  eight corner voxels of the cell containing the query point.
  \choice Bicubic interpolation has 16 unknowns, determined by the four corner intensities together
  with their first and mixed derivatives --- and since those derivatives must themselves be
  approximated from further pixels, bicubic effectively uses more than four samples.
  \choice Because bilinear interpolation is a weighted average of the surrounding values, applying
  it repeatedly during an iterative registration leaves the effective spatial resolution unchanged.
\end{choices}

\begin{solution}
  Answer: D.

  Averaging is exactly what costs you resolution. Bilinear interpolation smooths the image slightly
  every time it is applied, and during an optimisation loop it is applied thousands of times, so the
  blurring compounds. This is one reason the standard strategy is bilinear \emph{during} the
  optimisation, where speed matters, and bicubic for the \emph{final} resampling, where quality
  matters and the cost is paid once.
\end{solution}

% =============================================================== Question 8 ===
\question
Which of the following statements about the loss functions used for image restoration is
\textbf{incorrect}?

\begin{choices}
  \choice MAE, MSE, NMSE, PSNR and SSIM are all computed directly between the prediction and the
  ground-truth image.
  \choice SSIM compares many local patches of the two images and combines three components ---
  luminance, contrast and structure.
  \choice A perceptual loss passes both images through a previously trained network and compares
  intermediate features; that feature network must have been trained on the same task and the same
  modality as the restoration model, otherwise the comparison is meaningless.
  \choice An adversarial loss is a distributional distance defined over \emph{sets} of samples
  rather than between two paired images.
\end{choices}

\begin{solution}
  Answer: C.

  The feature network does not have to match the task or the modality. The lecture explicitly allows
  a well-established network trained on \emph{natural} images --- VGG is the standard choice --- or a
  network trained on CT images for some entirely different task. What is being borrowed is a
  general-purpose notion of what makes two images look different in terms of contextual content
  rather than pixel values, and that transfers surprisingly well across domains.

  The motivation for all of this is that MSE-type losses produce visibly over-smoothed output;
  combining a perceptual loss with an adversarial one gave the most realistic denoised low-dose CT
  in the referenced work.
\end{solution}

% =============================================================== Question 9 ===
\question
A semi-supervised auxiliary loss enforces consistency under spatial transformations. Two
transformations $\phi_1, \phi_2$ are drawn per image and the loss is
$\mathcal{L}\big(\phi_2 \circ \phi_1^{-1} \circ f(\phi_1 \circ x;\theta),\ f(\phi_2 \circ x;\theta)
\big)$. Which of the following statements is \textbf{incorrect}?

\begin{choices}
  \choice The loss requires no labels, so it can be evaluated on unlabelled images.
  \choice It encodes the principle that if $x \Leftrightarrow y$ then $\phi \circ x \Leftrightarrow
  \phi \circ y$, which is the same principle that underlies data augmentation.
  \choice The inverse $\phi_1^{-1}$ appears so that the two predictions are brought into a common
  spatial frame; without it they would sit in different frames and could not be compared pixel by
  pixel.
  \choice Because the loss is defined on unlabelled data, it replaces the supervised loss rather
  than being added to it.
\end{choices}

\begin{solution}
  Answer: D.

  It is an \emph{auxiliary} loss: the total cost is the supervised loss over the labelled examples
  \emph{plus} the consistency term, and the consistency term is summed over labelled and unlabelled
  examples alike. Dropping the supervised loss would remove the only thing that ties the network's
  output to the actual task --- a network that predicts a constant map everywhere is perfectly
  consistent under any transformation, and perfectly useless.
\end{solution}

% ============================================================== Question 10 ===
\question
When the network parameters are sampled from their posterior rather than optimised to a point
estimate, the sampled parameters can perform noticeably worse on the task than the MAP estimate.
What is the explanation given in the lecture?

\begin{choices}
  \choice In high dimensions the posterior is always unimodal, so every sample collapses onto the
  mean and no diversity is possible.
  \choice For an isotropic Gaussian in $d$ dimensions the expected distance of a sample from the
  mean grows like $\sqrt{d}$, so samples concentrate away from the mode --- and since the likelihood
  $p(D \mid \theta, M)$ \emph{is} the task loss, moving away from the mode means moving away from
  well-performing parameters.
  \choice High-dimensional posteriors cannot be sampled at all, so any apparent sample is an
  artefact of the random number generator.
  \choice In high dimensions the task loss becomes independent of the parameters, so performance is
  essentially random.
  \choice Sampling error decreases with dimension, so the effect disappears entirely for large
  networks.
\end{choices}

\begin{solution}
  Answer: B.

  This is the ``typical set'' phenomenon and it is thoroughly counter-intuitive: in high dimensions
  almost none of the probability mass sits near the mode, even though the mode is the single most
  probable point. Samples live in a thin shell at distance $\approx \sqrt{d}$ from the mean.

  The consequence for uncertainty estimation is uncomfortable. Because $\log p(D \mid \theta, M)$ is
  precisely the training objective, a parameter vector drawn honestly from the posterior can sit in
  a region of poor task performance --- so a faithful posterior sample and a well-performing network
  are not the same thing. Bounded sampling schemes have been proposed to mitigate this, but they do
  not fully solve it.
\end{solution}

% ============================================================== Question 11 ===
\question
A Vision Transformer is applied to images of size $224 \times 224$ with 3 channels. It uses a patch
size of $P = 16$, an embedding dimension of $D = 768$, and 12 attention heads. Recall that the image
is split into $N = HW/P^2$ patches, each flattened and mapped by a patch-embedding matrix
$E \in \mathbb{R}^{(P^2 C) \times D}$, and that a learnable class token is prepended to the sequence.
Circle True or False:

\begin{parts}
  \part \TFT\ \ The number of patch tokens is $N = 196$.

  \part \TFT\ \ The patch-embedding matrix $E$ is $768 \times 768$ and therefore contains
  $589\,824$ weights.

  \part \TFF\ \ After the class token is added, the sequence entering the encoder still has length
  196, because the class token takes the place of one patch.

  \part \TFF\ \ The attention matrix $QK^T$ of a single head has $N \times N$ entries, so its size
  grows linearly with the number of tokens.
\end{parts}

\begin{solution}
  (a) TRUE --- $N = HW/P^2 = (224 \times 224)/16^2 = 50\,176/256 = 196$, i.e.\ a $14 \times 14$ grid
  of patches.

  (b) TRUE --- the flattened patch has $P^2 C = 16^2 \times 3 = 256 \times 3 = 768$ entries, and it
  is mapped to $D = 768$, so $E$ is $768 \times 768 = 589\,824$. (That the two dimensions coincide is
  a property of this particular configuration, not a general rule.)

  (c) FALSE --- the class token is \emph{prepended}, not substituted: the sequence length becomes
  $196 + 1 = 197$. It carries no patch content of its own; it exists so that the final
  representation at that position can be read out and passed to the classification MLP.

  (d) FALSE --- an $N \times N$ matrix has $N^2$ entries, so the cost grows \textbf{quadratically}
  with sequence length. Here that is $197^2 = 38\,809$ attention scores per head, and 12 heads per
  layer. This quadratic scaling is the reason patching exists at all: attending over all
  $224 \times 224 = 50\,176$ pixels individually would need over $2.5$ billion scores per head.
\end{solution}

% ============================================================== Question 12 ===
\question
Which of the following statements on self-supervised learning are correct (multiple answers are
possible)?

\begin{choices}
  \choice Contrastive learning needs labelled examples in order to decide which pairs of images
  should be treated as positive and which as negative.
  \choice A pre-text task is a task that requires only images; the network trained on it is then
  fine-tuned with very few labelled examples on the task you actually care about.
  \choice In the \emph{local} contrastive loss the transformations must either preserve the patch
  correspondences or be undone before the corresponding patches are compared.
  \choice Foundation models in this setting are obtained by supervised training at very large scale
  on carefully curated labelled data sets.
  \choice The InfoNCE loss maximises the similarity between two views of the same image through its
  numerator, and minimises the similarity to the samples of the negative set through its
  denominator.
\end{choices}

\begin{solution}
  Answer: (B), (C) and (E) are correct.

  (A) is wrong, and it is wrong about the whole point of the method. The labels in contrastive
  learning are \emph{surrogate} labels manufactured from the data itself: two random transformations
  of the same image are declared ``close'', everything else ``far''. No human annotation is involved
  --- which is precisely why the approach is useful when labels are scarce. Note also that the
  transformations are not restricted to geometric ones; intensity changes are used deliberately, so
  that the representation becomes invariant to them.

  (D) inverts the defining property. Foundation models here are built by \textbf{self-supervised}
  training at very large scale --- masked auto-encoding and variations of contrastive learning ---
  on transformer architectures. Requiring curated labels at that scale is exactly what these methods
  exist to avoid.
\end{solution}

\end{questions}

\end{document}
